%%%%%%%%%%%%%%%%%%%%%%%%%%%%%%%%%%%%%%%%%%%%%%%%%%%%%%%%%%%%%%%%%%%%%%
% LaTeX Example: Project Report
%
% Source: http://www.howtotex.com
%
% Feel free to distribute this example, but please keep the referral
% to howtotex.com
% Date: March 2011 
% 
%%%%%%%%%%%%%%%%%%%%%%%%%%%%%%%%%%%%%%%%%%%%%%%%%%%%%%%%%%%%%%%%%%%%%%
% How to use writeLaTeX: 
%
% You edit the source code here on the left, and the preview on the
% right shows you the result within a few seconds.
%
% Bookmark this page and share the URL with your co-authors. They can
% edit at the same time!
%
% You can upload figures, bibliographies, custom classes and
% styles using the files menu.
%
% If you're new to LaTeX, the wikibook is a great place to start:
% http://en.wikibooks.org/wiki/LaTeX
%
%%%%%%%%%%%%%%%%%%%%%%%%%%%%%%%%%%%%%%%%%%%%%%%%%%%%%%%%%%%%%%%%%%%%%%
% Edit the title below to update the display in My Documents
%\title{Project Report}
%
%%% Preamble
\documentclass[paper=a4, fontsize=11pt]{scrartcl}
\usepackage[T1]{fontenc}
\usepackage{fourier}

\usepackage[english]{babel}															% English language/hyphenation
\usepackage[protrusion=true,expansion=true]{microtype}	
\usepackage{amsmath,amsfonts,amsthm} % Math packages
\usepackage[pdftex]{graphicx}	
\usepackage{url}
\usepackage{titlesec}
\usepackage{mathtools}
%\usepackage{lstlisting}


%%% Custom headers/footers (fancyhdr package)
\usepackage{fancyhdr}
\pagestyle{fancyplain}
\fancyhead{}											% No page header
\fancyfoot[L]{}											% Empty 
\fancyfoot[C]{}											% Empty
\fancyfoot[R]{\thepage}									% Pagenumbering
\renewcommand{\headrulewidth}{0pt}			% Remove header underlines
\renewcommand{\footrulewidth}{0pt}				% Remove footer underlines
\setlength{\headheight}{13.6pt}
\usepackage[margin=0.5in]{geometry}

%%% Equation and float numbering
\numberwithin{equation}{section}		% Equationnumbering: section.eq#
\numberwithin{figure}{section}			% Figurenumbering: section.fig#
\numberwithin{table}{section}				% Tablenumbering: section.tab#

%% Section Format
\titleformat*{\section}{\LARGE\bfseries}
\titleformat*{\subsection}{\Large\bfseries}
\titleformat*{\subsubsection}{\large\bfseries}
\titleformat*{\paragraph}{\large\bfseries}
\titleformat*{\subparagraph}{\large\bfseries}

\DeclarePairedDelimiter{\ceil}{\lceil}{\rceil}

%%% Maketitle metadata
\newcommand{\horrule}[1]{\rule{\linewidth}{#1}} 	% Horizontal rule

\title{
		%\vspace{-1in} 	
		\usefont{OT1}{bch}{b}{n}
		\normalfont \normalsize \textsc{University of Purdue\\
      Data Communication and Computer Networks\\
      CS 53600 Spring 2020} \\ [25pt]
		\horrule{0.5pt} \\[0.4cm]
		\huge Assignment 2\\
		\horrule{2pt} \\[0.5cm]
}
\author{
		\normalfont\normalsize Pedro da Costa Abreu Júnior\\ [-5pt]
    \normalsize Student ID 0030878383\\[-5pt]
    \normalsize \today
}
\date{}


%%% Begin document
\begin{document}
\maketitle

\section*{Problem 1}
HTTP, SMT, and POP3 uses TCP because they need to ensure that the data sent is delivered consistently, i.e. without errors.
DNS, on the other hand, prioritize performance and hence uses UDP.


\section*{Problem 2}
Let FTT be the file transmission time, and RTT be the roudtrip time, and D be the total delay time of the server
\subsection*{a}

$$FTT = \frac{\texttt{file size}}{R} = \frac{1 * 10^{3}b}{100 * 10^{6}b/s} = 10^{-5}s = 0.01ms$$
$$D = k * (2 * RTT + FTT) + 2 RTT = 100 * (2 * 1 + 10^{-5}) + 2 \approx 202s$$

\subsection*{b}
$$D = 2 * RTT + k * (RTT + FTT) = 2 * 1 + 100 * (1+10^{-5}) \approx 102s$$

\subsection*{c}
$$D = \ceil[\Bigg]{\frac{2 * RTT + k * FTT}{10}} \approx 20s$$

\subsection*{d}
As we can see with the formulas given above, c does not affect persistent connection, but may severely reduce the delay time for non-persistent connection.

Since R is very small it is dominated by the other variables and does not affect our calculations.

Non-persistent connection has to open a connection for each request, therefore the number of requests k will incur two extras RTT for the response. On the other hand, persistent connection does not need these two extras RTTs, meaning the value k has a much smaller effect on the final delay.

The bandwidth r will always have a direct effect on the communication, because it is the fundamental speed through which files will travel through the network.

\section*{Problem 3}
\subsection*{a}
using the command \verb|nslookup cs.purdue.edu| we get the Address 128.10.19.120.

\subsection*{b}
Using the command \verb|nslookup -query=ns purdue.edu| we get the following Authoritative Name Servers:

\begin{itemize}
\item ns2.purdue.edu.
\item pendragon.cs.purdue.edu.
\item ns.purdue.edu.
\item harbor.ecn.purdue.edu.
\end{itemize}

Using the command \verb|nslookup -query=ns cs.purdue.edu| we get the following Authoritative Name Servers:

\begin{itemize}
\item pendragon.cs.purdue.edu.
\item ns2.purdue.edu.
\item harbor.ecn.purdue.edu.
\item ns.purdue.edu.
\end{itemize}

\subsection*{c}
We can use the command \verb|nslookup -type=mx purdue.edu|
which yields the following results:

\begin{itemize}
\item xppmailspam01.itap.purdue.edu.
\item xppmailspam12.itap.purdue.edu.
\item xppmailspam02.itap.purdue.edu.
\item xppmailspam03.itap.purdue.edu.
\item xppmailspam05.itap.purdue.edu.
\item xppmailspam09.itap.purdue.edu.
\item xppmailspam08.itap.purdue.edu.
\item xppmailspam10.itap.purdue.edu.
\item xppmailspam11.itap.purdue.edu.
\item xppmailspam04.itap.purdue.edu.
\item xppmailspam06.itap.purdue.edu.
\item xppmailspam13.itap.purdue.edu.
\end{itemize}

\verb|nslookup -type=mx cs.purdue.edu| yields the following mail server mailslot.cs.purdue.edu.

%%% End document
\end{document}

%%% Local Variables:
%%% mode: latex
%%% TeX-master: t
%%% End:
